% GDZ PPN236018647,native563;HKUST564. Complete91a/b and commentary.
\clearpage\SourceLeaf{563}\ArithmeticHead{534}{L}{TAGEBUCH MIT ERLÄUTERUNGEN.}
\NumberedItem{91a}
\[
\begin{aligned}
&1+\frac19\frac{1\,.\,3}{4\,.\,4}+\frac1{81}\frac{1\,.\,3\,.\,5\,.\,7}{4\,.\,4\,.\,8\,.\,8}+\frac1{729}\frac{1\,.\,3\,.\,5\,.\,7\,.\,9\,.\,11}{4\,.\,4\,.\,8\,.\,8\,.\,12\,.\,12}\ldots=1{,}02220\ldots\\
&\qquad=\frac{1{,}3110\ldots}{3{,}1415\ldots}\surd6\left[=\frac\varpi2\frac1\pi\sqrt6\right].
\end{aligned}
\]
\par\noindent\hfill[1798] Iul.

\NumberedItem{91b}
\[
\operatorname{arc.\,sin\,lem.}\sin\varphi-\operatorname{arc.\,sin\,lem.}\cos\varphi=\varpi-\frac{2\varphi\varpi}\pi
\]
\[
\begin{aligned}
\operatorname{sin\,lemnisc.}[a]&=0{,}95500598\sin[a]\\
&\quad-0{,}0430495\ \sin3[a]\\
&\quad+0{,}0018605\ \sin5[a]\\
&\quad-0{,}0000803\ \sin7[a]
\end{aligned}
\]
\[
\begin{aligned}
\sin^2\operatorname{lem.}[a]&=0{,}4569472=\frac\pi{\varpi\varpi}\\
&\quad-[0{,}4569472]\cos2[a]\ldots
\end{aligned}
\]
\[
\begin{aligned}
\operatorname{arc.\,sin\,lem.}\sin\varphi
&=\frac\varpi\pi\varphi+\left(\frac\varpi\pi-\frac2\varpi\right)\sin2\varphi\\
&\quad+\left(\frac{11}2\frac\varpi\pi-\frac{12}\varpi\right)\sin4\varphi+\ldots
\end{aligned}
\]
\[
\begin{aligned}
\sin^5[\varphi]&=0{,}4775031\sin[\varphi]\\
&\quad+0{,}03\,.\,.\,.\,.\,.\,.\ [\sin3\varphi]\\
&\quad .\quad .\quad .\quad .\quad .\quad .
\end{aligned}
\]

\medskip\begingroup\small
\indent Die Aufzeichnungen der Nr. 91a und b stehen im engsten Zusammenhang mit den oben S. 168 ff. in den artt. [4.]—[6.] abgedruckten Notizen. Die Reihe der Nr. 91a findet sich im art. [5.], oben S. 169, der die Aufzeichnungen auf der Rückseite des Leistetitels wiedergibt, und dient dort zur Berechnung von \(\frac\varpi\pi\); \textsc{Gauss} hat sie später (im November 1799) auf S. 7 der Scheda Ac (abgedruckt oben S. 184 Gleichung [3]) noch einmal hingeschrieben.

Die trigonometrischen Reihen der Nr. 91b finden sich im wesentlichen im art. [4.], oben S. 168. Bei der Entwicklung von \(\operatorname{arc\,sin\,lem}s\) ist hier \(s=\sin\varphi\) gesetzt und \(\varphi\) im Bogenmaß gedacht, während S. 168 \(\varphi\) in Graden ausgedrückt erscheint; darum lautet das erste Glied rechts vom Gleichheitszeichen hier \(\frac\varpi\pi\varphi\), dort \(\varphi^0\). Für den Koeffizienten von \(\sin2\varphi\) sind auf S. 9 der Scheda Ac zwei, oben S. 184 Gln. [4], [5] abgedruckte Reihenentwicklungen angegeben. Die Entwicklung des art. [6.], oben S. 170, erscheint als unmittelbare Umkehrung der hier angegebenen Formel für \(\operatorname{arc\,sin\,lem}\sin\varphi\), wenn man dort \(\varphi\) mit \(\chi\) vertauscht, also \(\operatorname{sin\,lem}\chi=\operatorname{sin\,circ}\varphi\) setzt. In der Entwicklung von \(\operatorname{sin\,lem}a\) ist hier ebenso wie auf S. 168 auf \textls[120]{beiden} Seiten der Gleichung \(a\) im Gradmaß zu denken, man hätte also, behufs Zurückführung auf das Bogenmaß, für \(180^0\) links vom Gleichheitszeichen die halbe Lemniskatenlänge \(\varpi\), dagegen rechts vom Gleichheitszeichen die halbe Kreislänge \(\pi\) zu setzen. Vergl. die Bemerkung zu der folgenden Nr. 92.
\par\noindent\hfill\textsc{Klein. Schlesinger.}
\par\endgroup
